% --- Archon physics formalization source begin ---
% archon:physics
% archon:covers PhyXMiniProblems/problem_phyx_mini_0694.lean
% archon:source-report reports/phyx_mini/problem_phyx_mini_0694.source.json

\chapter{Physics problem phyx\_mini\_0694}
\label{ch:PhyXMiniProblems_problem_phyx_mini_0694}

\paragraph{Problem source.}
\#\# Physical scenario

The pincube machine was an ill-fated predecessor of the pinball machine. A \textbackslash{}(100\textbackslash{},\textbackslash{}mathrm\{g\}\textbackslash{}) cube is launched by pulling a spring back \textbackslash{}(20\textbackslash{},\textbackslash{}mathrm\{cm\}\textbackslash{}) and releasing it. The spring constant is \textbackslash{}(20\textbackslash{},\textbackslash{}mathrm\{N/m\}\textbackslash{}) and the surface is frictionless.

\#\# Question

What is the cube's launch speed, as it leaves the spring?

\#\# Answer choices

- A: A: 2.5m/s

- B: B: 3.1m/s

- C: C: 2.8m/s

- D: D: 3.4m/s

\#\# Figure caption (auxiliary; use the image as primary evidence)

The image depicts a horizontal spring-mass system. On the left, there's a compressed spring attached to a vertical wall. Beside the spring, there's a small block labeled with a mass, \textbackslash{}( m = 0.10 \textbackslash{}, \textbackslash{}text\{kg\} \textbackslash{}).

The block is displaced to the right, indicated by a horizontal arrow pointing in that direction. The displacement is labeled as \textbackslash{}( \textbackslash{}Delta r = 20 \textbackslash{}, \textbackslash{}text\{cm\} \textbackslash{}). 

There is also a vertical axis marked with \textbackslash{}( x \textbackslash{}) and a position labeled as the "Equilibrium position of spring," located to the right of the spring. There's a line marking the equilibrium position at \textbackslash{}( 0 \textbackslash{}).

The composition suggests a scenario where the spring is initially compressed by 20 cm from its equilibrium position.

\#\# Dataset metadata

Work and Energy; Physical Model Grounding Reasoning, Multi-Formula Reasoning

\paragraph{Recorded answer/context.}
C: C: 2.8m/s

\paragraph{Figure/image path.}
/root/proposal\_for\_physic/science-mango/phyx\_mini\_run/phyx\_data/test\_image/694.png

\paragraph{Formalization target.}
create a compiling Lean file with sorry bodies at `PhyXMiniProblems/problem\_phyx\_mini\_0694.lean`. The Lean declarations must preserve the physical quantities, dimensions or dimensional roles, figure labels, governing-law hypotheses, and final relation expressed by this problem.
Use Mathlib/Physlib names found through LeanExplore where available. If a physics API is missing, introduce faithful local abstractions rather than scalar placeholder aliases.

\begin{theorem}[Physics formalization target]
\label{thm:physics:phyx_mini_0694:target}
\lean{PhyXMiniProblems.ProblemPhyXMini0694.springCubeLaunchSpeed}
\uses{def:physics:phyx-mini-0694:phyxminiproblems-problemphyxmini0694-speedinmeterspersecond, def:physics:phyx-mini-0694:phyxminiproblems-problemphyxmini0694-springcubelaunchsetup, def:physics:phyx-mini-0694:phyxminiproblems-problemphyxmini0694-matchesfrictionlessspringlaunchscenario, def:physics:phyx-mini-0694:phyxminiproblems-problemphyxmini0694-matchesgivenproblemdata, def:physics:phyx-mini-0694:phyxminiproblems-problemphyxmini0694-matchesprimaryspringcubefigure, def:physics:phyx-mini-0694:phyxminiproblems-problemphyxmini0694-hasphysicalspringcubeparameters, def:physics:phyx-mini-0694:phyxminiproblems-problemphyxmini0694-satisfieselasticpotentialenergylaw, def:physics:phyx-mini-0694:phyxminiproblems-problemphyxmini0694-satisfiestranslationalkineticenergylaw, def:physics:phyx-mini-0694:phyxminiproblems-problemphyxmini0694-satisfiesmechanicalenergyconservation, def:physics:phyx-mini-0694:phyxminiproblems-problemphyxmini0694-agreeswithonedecimalspeedreadout, def:physics:phyx-mini-0694:phyxminiproblems-problemphyxmini0694-answerchoicereportslaunchspeed}
The assigned autoformalize agent should translate this physics problem into Lean declarations in the covered file, with theorem and lemma proof bodies written as `by sorry`.
\end{theorem}
\begin{proof}
This is an autoformalization task, not a proof task. Produce statements that can later be proved by the physics prover without weakening the physical model.
\end{proof}

% --- Archon declaration topology begin ---
\section{Declaration topology}
\begin{definition}[Length Quantity]
  \label{def:physics:phyx-mini-0694:phyxminiproblems-problemphyxmini0694-lengthquantity}
  \lean{PhyXMiniProblems.ProblemPhyXMini0694.LengthQuantity}
  A nonnegative, unit-independent physical length
\end{definition}
\begin{proof}
  This is the declared model definition of Length Quantity.
\end{proof}
\begin{definition}[Mass Quantity]
  \label{def:physics:phyx-mini-0694:phyxminiproblems-problemphyxmini0694-massquantity}
  \lean{PhyXMiniProblems.ProblemPhyXMini0694.MassQuantity}
  A nonnegative, unit-independent physical mass
\end{definition}
\begin{proof}
  This is the declared model definition of Mass Quantity.
\end{proof}
\begin{definition}[Spring Stiffness Quantity]
  \label{def:physics:phyx-mini-0694:phyxminiproblems-problemphyxmini0694-springstiffnessquantity}
  \lean{PhyXMiniProblems.ProblemPhyXMini0694.SpringStiffnessQuantity}
  A nonnegative, unit-independent spring stiffness. Since N/m = kg/s², its dimension is mass * time⁻²
\end{definition}
\begin{proof}
  This is the declared model definition of Spring Stiffness Quantity.
\end{proof}
\begin{definition}[length Readout]
  \label{def:physics:phyx-mini-0694:phyxminiproblems-problemphyxmini0694-lengthreadout}
  \lean{PhyXMiniProblems.ProblemPhyXMini0694.lengthReadout}
  \uses{def:physics:phyx-mini-0694:phyxminiproblems-problemphyxmini0694-lengthquantity}
  Read a physical length in a selected length unit
\end{definition}
\begin{proof}
  This is the declared model definition of length Readout.
\end{proof}
\begin{definition}[length In Meters]
  \label{def:physics:phyx-mini-0694:phyxminiproblems-problemphyxmini0694-lengthinmeters}
  \lean{PhyXMiniProblems.ProblemPhyXMini0694.lengthInMeters}
  \uses{def:physics:phyx-mini-0694:phyxminiproblems-problemphyxmini0694-lengthquantity, def:physics:phyx-mini-0694:phyxminiproblems-problemphyxmini0694-lengthreadout}
  Read a physical length in metres
\end{definition}
\begin{proof}
  This is the declared model definition of length In Meters.
\end{proof}
\begin{definition}[length In Centimeters]
  \label{def:physics:phyx-mini-0694:phyxminiproblems-problemphyxmini0694-lengthincentimeters}
  \lean{PhyXMiniProblems.ProblemPhyXMini0694.lengthInCentimeters}
  \uses{def:physics:phyx-mini-0694:phyxminiproblems-problemphyxmini0694-lengthquantity, def:physics:phyx-mini-0694:phyxminiproblems-problemphyxmini0694-lengthreadout}
  Read a physical length in centimetres
\end{definition}
\begin{proof}
  This is the declared model definition of length In Centimeters.
\end{proof}
\begin{definition}[mass Readout]
  \label{def:physics:phyx-mini-0694:phyxminiproblems-problemphyxmini0694-massreadout}
  \lean{PhyXMiniProblems.ProblemPhyXMini0694.massReadout}
  \uses{def:physics:phyx-mini-0694:phyxminiproblems-problemphyxmini0694-massquantity}
  Read a physical mass in a selected mass unit
\end{definition}
\begin{proof}
  This is the declared model definition of mass Readout.
\end{proof}
\begin{definition}[mass In Kilograms]
  \label{def:physics:phyx-mini-0694:phyxminiproblems-problemphyxmini0694-massinkilograms}
  \lean{PhyXMiniProblems.ProblemPhyXMini0694.massInKilograms}
  \uses{def:physics:phyx-mini-0694:phyxminiproblems-problemphyxmini0694-massquantity, def:physics:phyx-mini-0694:phyxminiproblems-problemphyxmini0694-massreadout}
  Read a physical mass in kilograms
\end{definition}
\begin{proof}
  This is the declared model definition of mass In Kilograms.
\end{proof}
\begin{definition}[mass In Grams]
  \label{def:physics:phyx-mini-0694:phyxminiproblems-problemphyxmini0694-massingrams}
  \lean{PhyXMiniProblems.ProblemPhyXMini0694.massInGrams}
  \uses{def:physics:phyx-mini-0694:phyxminiproblems-problemphyxmini0694-massquantity, def:physics:phyx-mini-0694:phyxminiproblems-problemphyxmini0694-massreadout}
  Read a physical mass in grams
\end{definition}
\begin{proof}
  This is the declared model definition of mass In Grams.
\end{proof}
\begin{definition}[speed In Meters Per Second]
  \label{def:physics:phyx-mini-0694:phyxminiproblems-problemphyxmini0694-speedinmeterspersecond}
  \lean{PhyXMiniProblems.ProblemPhyXMini0694.speedInMetersPerSecond}
  Read a speed in coherent SI units, metres per second
\end{definition}
\begin{proof}
  This is the declared model definition of speed In Meters Per Second.
\end{proof}
\begin{definition}[spring Stiffness In Newtons Per Meter]
  \label{def:physics:phyx-mini-0694:phyxminiproblems-problemphyxmini0694-springstiffnessinnewtonspermeter}
  \lean{PhyXMiniProblems.ProblemPhyXMini0694.springStiffnessInNewtonsPerMeter}
  \uses{def:physics:phyx-mini-0694:phyxminiproblems-problemphyxmini0694-springstiffnessquantity}
  Read a spring stiffness in coherent SI units, newtons per metre
\end{definition}
\begin{proof}
  This is the declared model definition of spring Stiffness In Newtons Per Meter.
\end{proof}
\begin{definition}[energy In Joules]
  \label{def:physics:phyx-mini-0694:phyxminiproblems-problemphyxmini0694-energyinjoules}
  \lean{PhyXMiniProblems.ProblemPhyXMini0694.energyInJoules}
  Read an energy in coherent SI units, joules
\end{definition}
\begin{proof}
  This is the declared model definition of energy In Joules.
\end{proof}
\begin{definition}[Launch Instant]
  \label{def:physics:phyx-mini-0694:phyxminiproblems-problemphyxmini0694-launchinstant}
  \lean{PhyXMiniProblems.ProblemPhyXMini0694.LaunchInstant}
  The two instants used in the energy balance
\end{definition}
\begin{proof}
  This is the declared model definition of Launch Instant.
\end{proof}
\begin{definition}[Spring Configuration]
  \label{def:physics:phyx-mini-0694:phyxminiproblems-problemphyxmini0694-springconfiguration}
  \lean{PhyXMiniProblems.ProblemPhyXMini0694.SpringConfiguration}
  The qualitative state of the spring at a launch instant
\end{definition}
\begin{proof}
  This is the declared model definition of Spring Configuration.
\end{proof}
\begin{definition}[Spring Mount]
  \label{def:physics:phyx-mini-0694:phyxminiproblems-problemphyxmini0694-springmount}
  \lean{PhyXMiniProblems.ProblemPhyXMini0694.SpringMount}
  How the spring is supported in the pictured apparatus
\end{definition}
\begin{proof}
  This is the declared model definition of Spring Mount.
\end{proof}
\begin{definition}[Surface Interaction]
  \label{def:physics:phyx-mini-0694:phyxminiproblems-problemphyxmini0694-surfaceinteraction}
  \lean{PhyXMiniProblems.ProblemPhyXMini0694.SurfaceInteraction}
  Frictional character of the horizontal contact surface
\end{definition}
\begin{proof}
  This is the declared model definition of Surface Interaction.
\end{proof}
\begin{definition}[Horizontal Direction]
  \label{def:physics:phyx-mini-0694:phyxminiproblems-problemphyxmini0694-horizontaldirection}
  \lean{PhyXMiniProblems.ProblemPhyXMini0694.HorizontalDirection}
  Horizontal directions used by the displacement and launch arrows
\end{definition}
\begin{proof}
  This is the declared model definition of Horizontal Direction.
\end{proof}
\begin{definition}[Axis Symbol]
  \label{def:physics:phyx-mini-0694:phyxminiproblems-problemphyxmini0694-axissymbol}
  \lean{PhyXMiniProblems.ProblemPhyXMini0694.AxisSymbol}
  Literal coordinate-axis symbols that a mechanics figure may display
\end{definition}
\begin{proof}
  This is the declared model definition of Axis Symbol.
\end{proof}
\begin{definition}[Spring Cube Figure]
  \label{def:physics:phyx-mini-0694:phyxminiproblems-problemphyxmini0694-springcubefigure}
  \lean{PhyXMiniProblems.ProblemPhyXMini0694.SpringCubeFigure}
  \uses{def:physics:phyx-mini-0694:phyxminiproblems-problemphyxmini0694-horizontaldirection, def:physics:phyx-mini-0694:phyxminiproblems-problemphyxmini0694-axissymbol}
  Raster-visible content of image 694.png. The two numerical fields are label readouts, not definitions of the corresponding physical quantities
\end{definition}
\begin{proof}
  This is the declared model definition of Spring Cube Figure.
\end{proof}
\begin{definition}[Spring Cube Launch Setup]
  \label{def:physics:phyx-mini-0694:phyxminiproblems-problemphyxmini0694-springcubelaunchsetup}
  \lean{PhyXMiniProblems.ProblemPhyXMini0694.SpringCubeLaunchSetup}
  \uses{def:physics:phyx-mini-0694:phyxminiproblems-problemphyxmini0694-lengthquantity, def:physics:phyx-mini-0694:phyxminiproblems-problemphyxmini0694-massquantity, def:physics:phyx-mini-0694:phyxminiproblems-problemphyxmini0694-springstiffnessquantity, def:physics:phyx-mini-0694:phyxminiproblems-problemphyxmini0694-launchinstant, def:physics:phyx-mini-0694:phyxminiproblems-problemphyxmini0694-springconfiguration, def:physics:phyx-mini-0694:phyxminiproblems-problemphyxmini0694-springmount, def:physics:phyx-mini-0694:phyxminiproblems-problemphyxmini0694-surfaceinteraction, def:physics:phyx-mini-0694:phyxminiproblems-problemphyxmini0694-horizontaldirection, def:physics:phyx-mini-0694:phyxminiproblems-problemphyxmini0694-springcubefigure}
  Independent physical quantities and state functions for the spring--cube system. In particular, the launch speed is an unconstrained physical field; it is not defined from an answer choice
\end{definition}
\begin{proof}
  This is the declared model definition of Spring Cube Launch Setup.
\end{proof}
\begin{definition}[Matches Frictionless Spring Launch Scenario]
  \label{def:physics:phyx-mini-0694:phyxminiproblems-problemphyxmini0694-matchesfrictionlessspringlaunchscenario}
  \lean{PhyXMiniProblems.ProblemPhyXMini0694.MatchesFrictionlessSpringLaunchScenario}
  \uses{def:physics:phyx-mini-0694:phyxminiproblems-problemphyxmini0694-lengthinmeters, def:physics:phyx-mini-0694:phyxminiproblems-problemphyxmini0694-speedinmeterspersecond, def:physics:phyx-mini-0694:phyxminiproblems-problemphyxmini0694-springcubelaunchsetup}
  The release-from-rest and equilibrium-departure conditions described by the problem. No field fixes the value of the departure speed
\end{definition}
\begin{proof}
  This is the declared model definition of Matches Frictionless Spring Launch Scenario.
\end{proof}
\begin{definition}[Matches Given Problem Data]
  \label{def:physics:phyx-mini-0694:phyxminiproblems-problemphyxmini0694-matchesgivenproblemdata}
  \lean{PhyXMiniProblems.ProblemPhyXMini0694.MatchesGivenProblemData}
  \uses{def:physics:phyx-mini-0694:phyxminiproblems-problemphyxmini0694-lengthincentimeters, def:physics:phyx-mini-0694:phyxminiproblems-problemphyxmini0694-massingrams, def:physics:phyx-mini-0694:phyxminiproblems-problemphyxmini0694-springstiffnessinnewtonspermeter, def:physics:phyx-mini-0694:phyxminiproblems-problemphyxmini0694-springcubelaunchsetup}
  Numerical data stated in the prose of the problem
\end{definition}
\begin{proof}
  This is the declared model definition of Matches Given Problem Data.
\end{proof}
\begin{definition}[Matches Primary Spring Cube Figure]
  \label{def:physics:phyx-mini-0694:phyxminiproblems-problemphyxmini0694-matchesprimaryspringcubefigure}
  \lean{PhyXMiniProblems.ProblemPhyXMini0694.MatchesPrimarySpringCubeFigure}
  \uses{def:physics:phyx-mini-0694:phyxminiproblems-problemphyxmini0694-lengthincentimeters, def:physics:phyx-mini-0694:phyxminiproblems-problemphyxmini0694-massinkilograms, def:physics:phyx-mini-0694:phyxminiproblems-problemphyxmini0694-springcubelaunchsetup}
  Literal evidence transcribed from the primary bitmap, including the displayed 0.10 kg, Δr = 20 cm, rightward arrow, and equilibrium marker at x = 0
\end{definition}
\begin{proof}
  This is the declared model definition of Matches Primary Spring Cube Figure.
\end{proof}
\begin{definition}[Has Physical Spring Cube Parameters]
  \label{def:physics:phyx-mini-0694:phyxminiproblems-problemphyxmini0694-hasphysicalspringcubeparameters}
  \lean{PhyXMiniProblems.ProblemPhyXMini0694.HasPhysicalSpringCubeParameters}
  \uses{def:physics:phyx-mini-0694:phyxminiproblems-problemphyxmini0694-lengthinmeters, def:physics:phyx-mini-0694:phyxminiproblems-problemphyxmini0694-massinkilograms, def:physics:phyx-mini-0694:phyxminiproblems-problemphyxmini0694-springstiffnessinnewtonspermeter, def:physics:phyx-mini-0694:phyxminiproblems-problemphyxmini0694-springcubelaunchsetup}
  Positivity conditions selecting a nondegenerate physical launch
\end{definition}
\begin{proof}
  This is the declared model definition of Has Physical Spring Cube Parameters.
\end{proof}
\begin{definition}[Satisfies Elastic Potential Energy Law]
  \label{def:physics:phyx-mini-0694:phyxminiproblems-problemphyxmini0694-satisfieselasticpotentialenergylaw}
  \lean{PhyXMiniProblems.ProblemPhyXMini0694.SatisfiesElasticPotentialEnergyLaw}
  \uses{def:physics:phyx-mini-0694:phyxminiproblems-problemphyxmini0694-lengthinmeters, def:physics:phyx-mini-0694:phyxminiproblems-problemphyxmini0694-springstiffnessinnewtonspermeter, def:physics:phyx-mini-0694:phyxminiproblems-problemphyxmini0694-energyinjoules, def:physics:phyx-mini-0694:phyxminiproblems-problemphyxmini0694-springcubelaunchsetup}
  Hooke-spring potential energy U = (1/2) k x² at both launch instants
\end{definition}
\begin{proof}
  This is the declared model definition of Satisfies Elastic Potential Energy Law.
\end{proof}
\begin{definition}[Satisfies Translational Kinetic Energy Law]
  \label{def:physics:phyx-mini-0694:phyxminiproblems-problemphyxmini0694-satisfiestranslationalkineticenergylaw}
  \lean{PhyXMiniProblems.ProblemPhyXMini0694.SatisfiesTranslationalKineticEnergyLaw}
  \uses{def:physics:phyx-mini-0694:phyxminiproblems-problemphyxmini0694-massinkilograms, def:physics:phyx-mini-0694:phyxminiproblems-problemphyxmini0694-speedinmeterspersecond, def:physics:phyx-mini-0694:phyxminiproblems-problemphyxmini0694-energyinjoules, def:physics:phyx-mini-0694:phyxminiproblems-problemphyxmini0694-springcubelaunchsetup}
  Translational kinetic energy K = (1/2) m v² at both launch instants
\end{definition}
\begin{proof}
  This is the declared model definition of Satisfies Translational Kinetic Energy Law.
\end{proof}
\begin{definition}[total Mechanical Energy In Joules]
  \label{def:physics:phyx-mini-0694:phyxminiproblems-problemphyxmini0694-totalmechanicalenergyinjoules}
  \lean{PhyXMiniProblems.ProblemPhyXMini0694.totalMechanicalEnergyInJoules}
  \uses{def:physics:phyx-mini-0694:phyxminiproblems-problemphyxmini0694-energyinjoules, def:physics:phyx-mini-0694:phyxminiproblems-problemphyxmini0694-launchinstant, def:physics:phyx-mini-0694:phyxminiproblems-problemphyxmini0694-springcubelaunchsetup}
  Total mechanical energy of the spring--cube system at one instant
\end{definition}
\begin{proof}
  This is the declared model definition of total Mechanical Energy In Joules.
\end{proof}
\begin{definition}[Satisfies Mechanical Energy Conservation]
  \label{def:physics:phyx-mini-0694:phyxminiproblems-problemphyxmini0694-satisfiesmechanicalenergyconservation}
  \lean{PhyXMiniProblems.ProblemPhyXMini0694.SatisfiesMechanicalEnergyConservation}
  \uses{def:physics:phyx-mini-0694:phyxminiproblems-problemphyxmini0694-springcubelaunchsetup, def:physics:phyx-mini-0694:phyxminiproblems-problemphyxmini0694-totalmechanicalenergyinjoules}
  With no dissipative work, total mechanical energy is conserved
\end{definition}
\begin{proof}
  This is the declared model definition of Satisfies Mechanical Energy Conservation.
\end{proof}
\begin{definition}[Answer Choice]
  \label{def:physics:phyx-mini-0694:phyxminiproblems-problemphyxmini0694-answerchoice}
  \lean{PhyXMiniProblems.ProblemPhyXMini0694.AnswerChoice}
  The four answer labels printed by the problem
\end{definition}
\begin{proof}
  This is the declared model definition of Answer Choice.
\end{proof}
\begin{definition}[answer Choice Speed In Meters Per Second]
  \label{def:physics:phyx-mini-0694:phyxminiproblems-problemphyxmini0694-answerchoicespeedinmeterspersecond}
  \lean{PhyXMiniProblems.ProblemPhyXMini0694.answerChoiceSpeedInMetersPerSecond}
  \uses{def:physics:phyx-mini-0694:phyxminiproblems-problemphyxmini0694-answerchoice}
  The speed in metres per second printed next to an answer label
\end{definition}
\begin{proof}
  This is the declared model definition of answer Choice Speed In Meters Per Second.
\end{proof}
\begin{definition}[Agrees With One Decimal Speed Readout]
  \label{def:physics:phyx-mini-0694:phyxminiproblems-problemphyxmini0694-agreeswithonedecimalspeedreadout}
  \lean{PhyXMiniProblems.ProblemPhyXMini0694.AgreesWithOneDecimalSpeedReadout}
  \uses{def:physics:phyx-mini-0694:phyxminiproblems-problemphyxmini0694-speedinmeterspersecond}
  Agreement with a speed displayed to the nearest tenth of a metre per second
\end{definition}
\begin{proof}
  This is the declared model definition of Agrees With One Decimal Speed Readout.
\end{proof}
\begin{definition}[Answer Choice Reports Launch Speed]
  \label{def:physics:phyx-mini-0694:phyxminiproblems-problemphyxmini0694-answerchoicereportslaunchspeed}
  \lean{PhyXMiniProblems.ProblemPhyXMini0694.AnswerChoiceReportsLaunchSpeed}
  \uses{def:physics:phyx-mini-0694:phyxminiproblems-problemphyxmini0694-springcubelaunchsetup, def:physics:phyx-mini-0694:phyxminiproblems-problemphyxmini0694-answerchoice, def:physics:phyx-mini-0694:phyxminiproblems-problemphyxmini0694-answerchoicespeedinmeterspersecond, def:physics:phyx-mini-0694:phyxminiproblems-problemphyxmini0694-agreeswithonedecimalspeedreadout}
  An answer choice reports the launch speed to the displayed precision
\end{definition}
\begin{proof}
  This is the declared model definition of Answer Choice Reports Launch Speed.
\end{proof}
% --- Archon declaration topology end ---
% --- Archon physics formalization source end ---
